\documentclass[a4paper]{article}
\usepackage[pdftex]{graphicx}
\usepackage[utf8]{inputenc}
\usepackage{enumerate}
\usepackage{amssymb}
\usepackage{eurosym}
\usepackage{tikz}
\usepackage{href-ul}
\hypersetup{
	colorlinks=true,
	linkcolor=blue,
	urlcolor=blue}
\usepackage{geometry}
\geometry{a4paper, top=15mm, left=15mm, right=15mm, bottom=15mm,
	headsep=10mm, footskip=12mm}

\begin{document}
\begin{tikzpicture}(10,3)
	\draw[ultra thick](2,0) --(10,0) -- (10,1.3) --(2,1.3) -- (2,0);
	\draw[fill=black](2,0)-- (0,.6) -- (2,.6) -- (2,0);
	\node at (6,.6) {\large Lösungsblatt von \href{https://www.okuyakl.de}{www.okuyakl.de}};
\end{tikzpicture}
\vspace{0.5 cm}

\noindent{\bf Aufgabe 1. a)}\\
Für $E_1$="Augensumme 2" gibt es folgende Möglichkeiten:
$$E_1=\{20;02;11\}$$
Die Wahrscheinlichkeit, dass ein Würfel eine 2 zeigt, ist $ 1 \over 6$, für eine 1 ist sie $ 1\over 2$ und für eine 0 ist sie $ 1\over 3$. Also ist die Wahrscheinlichkeit für die Augensumme 2:
$$p(E_1)= 2 \cdot   {1 \over 6} \cdot  {1 \over 3} +   \left({1 \over 2}\right)^2={13\over 36}$$

\noindent{\bf Aufgabe 1. b)}\\
$$
\renewcommand{\arraystretch}{2}
	\begin{array}{c|c|c|c|c|c}
		 x_i  & 0 & 1 & 2 & 3 & 4 \\
		\hline
		 P(X=x_i)  & {1\over9 } & {1\over 3}  &  \frac{13}{36}   & {1\over 6}   & {1\over 36}\\
	\end{array}
$$

\noindent{\bf Aufgabe 1. c)}\\
Für $E_2$="3 gleichfarbige Kugeln" gibt es folgende Möglichkeiten:
$$E_2=\{rrr;ggg;bbb\}$$
Die Wahrscheinlichkeit hierfür ist:
$$p(E_2)= 3 \cdot {3 \over 9} \cdot {2\over8} \cdot {1\over 7}= {1\over 28}$$

\noindent{\bf Aufgabe 1. d)}\\
Hierfür stellen wir einen Term für den Erwartungswert  aus Sicht der Sendeleitung auf :
$$
\renewcommand{\arraystretch}{2}
\begin{array}{rcll}
E(x)&=& (2-B) \cdot p(E_2)+ 2 \cdot (1-p(E_2))\\
1,25 &=& {2\over 28}- B\cdot {1\over 28} +{54\over 28}\\
B &=& 0,75 \cdot 28 = 21\\
\end{array}
$$
Der Gewinn muss \EUR{21} betragen.

\noindent{\bf Aufgabe 2. a)}\\

Die Wahrscheinlichkeit für das Ereignis (i) entspricht einer Bernoullikette der Länge 10:

$$P_{1 \over 3}^{10}(X=5)=B(10;{1 \over 3}; 5)={10 \choose 5} \cdot  \left( {1 \over 3} \right)^5 \cdot  
\left( {2 \over 3} \right)^5=13,6\%$$

Für die Wahrscheinlichkeit des Ereignisses (ii) formen wir um und lesen aus dem Tafelwerk ab:

$$ P_{1 \over 3}^{10}(4\le X \le 8)= F_{1 \over 3}^{10}(8)-F_{1 \over 3}^{10}(3)=0,999644-0,559264=44,0\%$$

\noindent {\bf Aufgabe 2. b)}

Wir erstellen folgende Tabelle der Wahrscheinlichkeitsverteilung:

$$
\begin{array}{c|c|c|c}
X_i & -2,90 & 0,10 & 1,10 \\ 
\hline
P(X_i) &  {1 \over 3} &  {1 \over 6} &  {1 \over 2}\\  
\end{array}
$$

Der Erwartungswert $E(x)$ der Zufallsgröße ,,Gewinn" beträgt:

$$E(x)=  -2,90 \cdot {1 \over 3} + 0,1 \cdot  {1 \over 6} + 1,1 \cdot  {1 \over 2}=-0,4$$

Man macht also durchschnittlich \EUR{0,4} Verlust pro Spiel.

\noindent {\bf Aufgabe 3. a)}

Die gesuchte Wahrscheinlichkeit ist wieder eine Bernoullikette:

$$P_{0,03}^{12}(X=2)=B(12;0,03; 2)={12 \choose 2} \cdot  0,03^2 \cdot  0,97^{10}=4,4\%$$

\newpage
\noindent {\bf Aufgabe 3. b)}

Dies ist eine sogenannte Dreimal-Mindestens-Aufgabe: Wir arbeiten mit dem Gegenereignis:

$$
\begin{array}{rcl}
1-P_{p}^{12}(X=0) & \ge & 0,9\\
1-(1-p)^{12}& \ge & 0,9\\
(1-p)^{12} & \le & 0,1  \quad | \sqrt[12]{}\\
1-p & \le & 0,8254 \\
p & \ge & 0,18 \quad \textnormal{also} \quad p \ge 18\%
\end{array}
$$

\noindent {\bf Aufgabe 4.}
Die Wahrscheinlichkeitsverteilung für $Y$
ist: 
$$
\renewcommand{\arraystretch}{2}
\begin{array}{|c|c|c|c|}
\hline
y_i & 1450 & 1900 & 2350 \\
\hline
P(Y=y_i)& 0,6 & 0,1 & 0,3 \\
\hline
\end{array}
$$
Der Erwartungswert hierbei ist:
$$E(Y)=1450\cdot 0,6 + 1900 \cdot 0,1 +2350 \cdot 0,3 = 1765$$
Die Varianz ist:
$$Var(Y)=(1450-1765)^2\cdot 0,6 +(1900-1765)^2\cdot 0,1+(2350-1765)^2\cdot 0,3 =164025$$
Damit ist die Standardabweichung:
$$\sigma(Y)=\sqrt{164025}=405$$

\begin{center}
	\includegraphics[width=7 cm]{../../viecher/endcomic.pdf}
	
	Hier geht es zurück zum \href{https://www.okuyakl.de/math/m12stoL064/aa064.pdf}{Aufgabenblatt}
\end{center}

\end{document}

